\section{Scelte progettuali e confine con la specifica}
Il modello corrente privilegia poche entità con codifiche compatte: \code{Request} per due tipi di pratica e \code{Sale} per vendita, acquisizione e penale. Questo semplifica le tabelle e i DAO, ma richiede controlli applicativi sul segno e sui riferimenti nulli. La Strategy separa gli algoritmi di prezzo dai servizi; l'Observer separa i subscriber dai produttori di eventi, pur senza garanzia after-commit. La documentazione precedente propone invece \code{Payment}, \code{PurchaseProposal}, \code{SaleOrder}, dashboard e stock order: sono scelte progettuali di un'altra baseline, non componenti del checkout illustrato nei diagrammi nuovi.
